\documentclass[12pt,a4paper]{article}

\usepackage[utf8]{inputenc}
\usepackage[T1]{fontenc}
\usepackage{amsmath,amssymb,amsthm}
\usepackage{mathtools}
\usepackage{tikz-cd}
\usepackage{listings}
\usepackage{xcolor}
\usepackage{geometry}
\usepackage{hyperref}
\usepackage{enumitem}
\usepackage{tcolorbox}
\usepackage{booktabs}
\usepackage{fancyhdr}

\geometry{margin=2.5cm}
\hypersetup{colorlinks=true, linkcolor=blue!60!black, urlcolor=blue!60!black}

% Theorem environments
\theoremstyle{definition}
\newtheorem{definition}{Definition}[section]
\newtheorem{example}[definition]{Example}

\theoremstyle{plain}
\newtheorem{theorem}[definition]{Theorem}
\newtheorem{proposition}[definition]{Proposition}

\theoremstyle{remark}
\newtheorem*{remark}{Remark}
\newtheorem*{reallife}{Real-Life Analogy}

% Code listing style
\definecolor{codebg}{RGB}{245,245,250}
\definecolor{codegreen}{RGB}{0,128,0}
\definecolor{codepurple}{RGB}{128,0,128}
\definecolor{codeblue}{RGB}{0,0,180}

\lstdefinelanguage{Dart}{
  keywords={class, sealed, abstract, final, static, void, int, String, bool,
            return, if, else, switch, case, Future, async, await, typedef,
            extension, on, List, Map, enum, get, required, true, false, null},
  keywordstyle=\color{codeblue}\bfseries,
  commentstyle=\color{codegreen}\itshape,
  stringstyle=\color{codepurple},
  morecomment=[l]{//},
  morecomment=[s]{/*}{*/},
  morestring=[b]',
  morestring=[b]",
  sensitive=true,
}

\lstset{
  language=Dart,
  backgroundcolor=\color{codebg},
  basicstyle=\ttfamily\small,
  breaklines=true,
  frame=single,
  rulecolor=\color{black!20},
  numbers=left,
  numberstyle=\tiny\color{gray},
  tabsize=2,
  xleftmargin=1.5em,
  framexleftmargin=1em,
}

% Custom box for "you already know this"
\newtcolorbox{youknow}[1][]{
  colback=green!5!white,
  colframe=green!50!black,
  title={\textbf{You Already Know This}},
  fonttitle=\sffamily,
  #1
}

% Custom box for the formal definition
\newtcolorbox{formal}[1][]{
  colback=blue!3!white,
  colframe=blue!40!black,
  title={\textbf{Formal Definition}},
  fonttitle=\sffamily,
  #1
}

\title{\Huge\bfseries The Category Theory You Already Know\\[0.5em]
  \Large A Programmer's Guide, Illustrated with a Real Codebase}
\author{Generated from categorical analysis of the \emph{Tech World} game\\[0.3em]
  \small A Flutter/Dart/Flame multiplayer educational game}
\date{May 2026}

\pagestyle{fancy}
\fancyhf{}
\fancyhead[L]{\small\itshape The Category Theory You Already Know}
\fancyhead[R]{\small\thepage}
\renewcommand{\headrulewidth}{0.4pt}

\begin{document}

\maketitle

\begin{abstract}
Category theory has a reputation for being abstract and impenetrable.
But if you write code, you already use categorical structures every day
--- you just don't call them that.
This document catalogues the category theory found hiding inside a real
production codebase: \emph{Tech World}, a multiplayer educational game
built with Flutter, Dart, and Flame.
Every definition is paired with (a) the actual code that implements it,
and (b) a real-life analogy.
No prior knowledge of category theory is assumed.
\end{abstract}

\tableofcontents
\newpage

%=============================================================================
\section{What Is a Category?}
%=============================================================================

Before we look at the code, let's nail down the one definition everything
else builds on.

\begin{formal}
A \textbf{category} $\mathcal{C}$ consists of:
\begin{enumerate}
  \item A collection of \textbf{objects}: $A, B, C, \ldots$
  \item For every pair of objects $A, B$, a collection of \textbf{morphisms}
        (arrows): $f \colon A \to B$
  \item A \textbf{composition} operation: given $f \colon A \to B$ and
        $g \colon B \to C$, there exists $g \circ f \colon A \to C$
  \item An \textbf{identity} morphism $\mathrm{id}_A \colon A \to A$ for every object $A$
\end{enumerate}
Subject to two laws:
\begin{itemize}
  \item \textbf{Associativity}: $h \circ (g \circ f) = (h \circ g) \circ f$
  \item \textbf{Identity}: $f \circ \mathrm{id}_A = f = \mathrm{id}_B \circ f$
\end{itemize}
\end{formal}

\begin{reallife}
Think of a city's road network.
Objects are intersections.
Morphisms are routes between intersections.
Composition is ``take route $f$, then take route $g$.''
The identity morphism is ``stay where you are.''
Associativity says it doesn't matter how you group a multi-leg journey
--- the overall route is the same.
\end{reallife}

\subsection{The Types Category: \textbf{Dart}}

The most important category for a programmer is the one you live in every day.

\begin{youknow}
In Dart (or any typed language), you are working in a category:
\begin{itemize}
  \item \textbf{Objects} = Dart types (\texttt{int}, \texttt{String},
        \texttt{BotStatus}, \texttt{AppEvent}, \ldots)
  \item \textbf{Morphisms} = functions between types
        (\texttt{int Function(String)}, etc.)
  \item \textbf{Composition} = function composition:
        \texttt{(f, g) $\mapsto$ (x) => g(f(x))}
  \item \textbf{Identity} = the identity function: \texttt{(x) => x}
\end{itemize}
We call this category $\mathbf{Dart}$.
\end{youknow}

Every time you write a function, you are defining a morphism. Every time you
chain two function calls, you are composing morphisms. You have been doing
category theory since your first ``Hello, World.''


\subsection{The State Category}

Games have states --- and transitions between states form a category too.

\begin{formal}
A \textbf{state machine} defines a category where:
\begin{itemize}
  \item Objects = states
  \item Morphisms = valid transitions between states
  \item Composition = sequencing transitions
  \item Identity = ``stay in the current state''
\end{itemize}
\end{formal}

In Tech World, the LiveKit connection has an explicit state machine:

\begin{lstlisting}
enum _ConnectionState {
  disconnected,  // not connected to any room
  connecting,    // handshake in progress
  connected,     // fully connected, media flowing
}
\end{lstlisting}

The transitions are monotone along the order
$\texttt{disconnected} < \texttt{connecting} < \texttt{connected}$:
you must pass through \texttt{connecting} to reach \texttt{connected}.
Reconnection resets to \texttt{disconnected} and climbs again.

\begin{reallife}
A traffic light: Red $\to$ Green $\to$ Amber $\to$ Red.
Each colour is an object; each legal change is a morphism.
You can't jump from Red to Amber --- the transitions are constrained,
just like the connection state machine above.
\end{reallife}


%=============================================================================
\section{Coproducts: Sum Types You Use Every Day}
\label{sec:coproducts}
%=============================================================================

\begin{formal}
A \textbf{coproduct} of objects $A$ and $B$ in a category $\mathcal{C}$ is
an object $A + B$ together with injection morphisms
$\iota_A \colon A \to A + B$ and $\iota_B \colon B \to A + B$
such that for any object $C$ with morphisms $f \colon A \to C$ and
$g \colon B \to C$, there exists a unique morphism
$[f, g] \colon A + B \to C$ making the diagram commute:
\[
\begin{tikzcd}
A \arrow[r, "\iota_A"] \arrow[dr, "f"'] & A + B \arrow[d, dashed, "{[f,g]}"] & B \arrow[l, "\iota_B"'] \arrow[dl, "g"] \\
& C &
\end{tikzcd}
\]
The unique morphism $[f,g]$ is called the \textbf{copairing}.
\end{formal}

\begin{youknow}
In Dart 3, a \texttt{sealed class} is exactly a coproduct.
Tech World's event system is a 34-way coproduct:
\begin{lstlisting}
sealed class AppEvent {
  PiiPolicy get piiPolicy;  // every summand must provide this
  Map<String, dynamic> toJson();
}

class RoomJoined extends AppEvent { ... }
class ChatSent extends AppEvent { ... }
class SpellCast extends AppEvent { ... }
// ... 31 more subclasses
\end{lstlisting}
Each subclass is an injection $\iota_i \colon A_i \to \texttt{AppEvent}$.
An exhaustive \texttt{switch} statement is the copairing $[f_1, \ldots, f_{34}]$
--- the compiler \emph{forces} you to handle every case, which is exactly the
universal property of the coproduct!
\end{youknow}

Here is another beautiful example. Tech World has two kinds of challenges
--- code challenges and prompt challenges --- unified into a single type:

\begin{lstlisting}
sealed class ChallengeRef {
  String get wireName;

  static ChallengeRef? parse(String wire) {
    // Try CodeChallengeId first, then PromptChallengeId
    final code = CodeChallengeId.parse(wire);
    if (code != null) return CodeRef(code);
    final prompt = PromptChallengeId.parse(wire);
    if (prompt != null) return PromptRef(prompt);
    return null;
  }
}

class CodeRef extends ChallengeRef { ... }
class PromptRef extends ChallengeRef { ... }
\end{lstlisting}

The \texttt{parse} function is a copairing: it takes a wire string and injects
it into the correct summand of the coproduct
$\texttt{ChallengeRef} = \texttt{CodeRef} + \texttt{PromptRef}$.

\begin{reallife}
A hospital emergency room triages patients into categories:
trauma, cardiac, respiratory, etc.
Each category is a summand.
The triage nurse is the copairing --- she examines each patient (the input)
and sends them to exactly one department.
The hospital's guarantee that \emph{every} patient gets triaged
is the universal property.
\end{reallife}

\subsection{The PiiPolicy Poset: A Two-Element Category}

Not all categories are big. Tech World has a category with just two objects:

\begin{lstlisting}
enum PiiPolicy { none, pii }
\end{lstlisting}

This is a \textbf{poset category} (partially ordered set viewed as a category):
$\texttt{none} \leq \texttt{pii}$, where there is a unique morphism
$\texttt{none} \to \texttt{pii}$ (``escalate to PII'') and identity morphisms.
Every \texttt{AppEvent} carries a \texttt{piiPolicy} field, making
\texttt{piiPolicy} a functor from the event category to this two-element poset.

The event dispatch system uses this: remote sinks (which send data over the network)
filter out events where $\texttt{piiPolicy} = \texttt{pii}$.
This filter is a \textbf{monotone map} --- it respects the ordering, ensuring
that PII never leaks to remote observers.


%=============================================================================
\section{Functors: Structure-Preserving Maps}
\label{sec:functors}
%=============================================================================

\begin{formal}
A \textbf{functor} $F \colon \mathcal{C} \to \mathcal{D}$ between categories
assigns:
\begin{itemize}
  \item To each object $A$ in $\mathcal{C}$, an object $F(A)$ in $\mathcal{D}$
  \item To each morphism $f \colon A \to B$, a morphism
        $F(f) \colon F(A) \to F(B)$
\end{itemize}
Preserving:
\begin{itemize}
  \item \textbf{Identity}: $F(\mathrm{id}_A) = \mathrm{id}_{F(A)}$
  \item \textbf{Composition}: $F(g \circ f) = F(g) \circ F(f)$
\end{itemize}
\end{formal}

\begin{youknow}
If you have ever written \texttt{list.map(f)}, you have used a functor.
\texttt{List} is a functor from $\mathbf{Dart}$ to $\mathbf{Dart}$:
it sends each type \texttt{T} to the type \texttt{List<T>},
and each function \texttt{f: A -> B} to the function
\texttt{List<A> -> List<B>} given by \texttt{.map(f)}.
\end{youknow}

\subsection{The Stream Functor}

Tech World uses Dart \texttt{Stream}s extensively for real-time data from LiveKit:

\begin{lstlisting}
// In LiveKit service: transforming raw data into typed events
_subscription = dataStream
    .map((bytes) => utf8.decode(bytes))   // Stream<Uint8List> -> Stream<String>
    .map((s) => jsonDecode(s))            // Stream<String> -> Stream<dynamic>
    .where((json) => json['topic'] != null)
    .map((json) => parseEvent(json));     // Stream<dynamic> -> Stream<GameEvent>
\end{lstlisting}

\texttt{Stream} is a functor: it maps types to streams-of-that-type, and
\texttt{.map(f)} lifts a function on elements to a function on streams.
The functor laws guarantee that:
\begin{align}
\texttt{stream.map((x) => x)} &= \texttt{stream} \tag{Identity} \\
\texttt{stream.map(f).map(g)} &= \texttt{stream.map((x) => g(f(x)))} \tag{Composition}
\end{align}

\begin{reallife}
A factory assembly line is a functor.
Raw materials (objects) become products.
Each processing step (morphism) transforms materials.
\texttt{.map} is ``apply this processing step to every item on the conveyor belt.''
The identity law says ``a do-nothing step changes nothing.''
The composition law says ``two steps in sequence equals one combined step''
--- rearranging the factory floor doesn't change the output.
\end{reallife}

\subsection{The Direction Functor: Enums to Vectors}

Here's a less obvious functor. Tech World maps compass directions to 2D vectors:

\begin{lstlisting}
enum Direction {
  up(0, -1),
  down(0, 1),
  left(-1, 0),
  right(1, 0);

  final int dx, dy;
  const Direction(this.dx, this.dy);
}
\end{lstlisting}

This is a functor from the discrete category of directions (four objects,
only identity morphisms) to $\mathbb{R}^2$.
It maps each direction to its offset vector: $\texttt{up} \mapsto (0, -1)$,
$\texttt{right} \mapsto (1, 0)$, etc.
Since the source category is discrete (no non-identity morphisms),
the functor laws are trivially satisfied.

\begin{reallife}
A dictionary is a functor from a discrete category of words to meanings.
Each word maps to its definition; there are no ``morphisms between words''
to worry about.
\end{reallife}


%=============================================================================
\section{Natural Transformations: Conversions That Don't Peek Inside}
\label{sec:natural-transformations}
%=============================================================================

\begin{formal}
Given functors $F, G \colon \mathcal{C} \to \mathcal{D}$, a
\textbf{natural transformation} $\alpha \colon F \Rightarrow G$
assigns to each object $A$ in $\mathcal{C}$ a morphism
$\alpha_A \colon F(A) \to G(A)$ in $\mathcal{D}$, such that for every
morphism $f \colon A \to B$ in $\mathcal{C}$, the following
\textbf{naturality square} commutes:
\[
\begin{tikzcd}
F(A) \arrow[r, "\alpha_A"] \arrow[d, "F(f)"'] & G(A) \arrow[d, "G(f)"] \\
F(B) \arrow[r, "\alpha_B"'] & G(B)
\end{tikzcd}
\]
That is: $\alpha_B \circ F(f) = G(f) \circ \alpha_A$.
In words: ``it doesn't matter whether you transform first and then convert,
or convert first and then transform.''
\end{formal}

\begin{youknow}
Tech World's \texttt{wireName}/\texttt{parse} pattern on enums is a
natural transformation.
Every domain enum (\texttt{CodeChallengeId}, \texttt{PromptChallengeId},
\texttt{WordId}, \texttt{LiveKitTopic}, \ldots) has:
\begin{lstlisting}
final String wireName;           // domain -> wire
static T? parse(String wire);   // wire -> domain (partial inverse)
\end{lstlisting}
\texttt{wireName} converts domain values to wire strings \emph{without
looking at what the value means} --- it's a uniform conversion that
works the same way for every enum. That uniformity is naturality.
\end{youknow}

\subsection{Serialisation as a Natural Transformation}

The JSON serialisation in Tech World is another natural transformation:
\begin{lstlisting}
// On every AppEvent subclass:
Map<String, dynamic> toJson();

// On MapEditBatch:
factory MapEditBatch.fromJson(Map<String, dynamic> json);
\end{lstlisting}

Let $F$ be the identity functor on domain types and $G$ be the
``JSON representation'' functor. Then \texttt{toJson} is a natural
transformation $\alpha \colon F \Rightarrow G$, and \texttt{fromJson}
is an approximate inverse $\beta \colon G \Rightarrow F$.

The naturality condition says: if you transform a domain object
(say, modify a map edit) and \emph{then} serialise, you get the same
JSON as if you serialise first and then transform the JSON.
This is exactly what makes the LiveKit data channel work ---
map edits serialised on one client deserialise correctly on another.

\begin{reallife}
Translation between languages is a natural transformation.
If you have a recipe (an object) and you want to double the quantities
(a morphism), it doesn't matter whether you:
\begin{enumerate}
  \item Double the quantities, then translate to French; or
  \item Translate to French, then double the quantities.
\end{enumerate}
You get the same French recipe either way.
The translation is ``natural'' because it doesn't depend on the
specific quantities --- it works uniformly.
\end{reallife}

\subsection{The Logger Bridge: A Filtering Natural Transformation}

Here is a subtler example. Tech World bridges Dart's logging system into
its event system:

\begin{lstlisting}
// lib/events/logger_bridge.dart
AppLogRecord? mapLogRecord(LogRecord record) {
  if (record.level < Level.INFO) return null;  // drop FINE
  return AppLogRecord(
    level: record.level.name,
    message: record.message,
    // ...
  );
}
\end{lstlisting}

This is a \emph{partial} natural transformation from the Logger functor to
the EventSink functor. The partiality (returning \texttt{null} for
FINE-level logs) is deliberate: it ensures that verbose debug logs
--- which might contain PII --- never reach the event dispatch system.
This is a privacy gate implemented as a mathematical structure.


%=============================================================================
\section{Monads: Composing Effectful Computations}
\label{sec:monads}
%=============================================================================

Monads are the most famous (and most feared) concept from category theory
in programming. But you've been using them all along.

\begin{formal}
A \textbf{monad} on a category $\mathcal{C}$ is a functor
$T \colon \mathcal{C} \to \mathcal{C}$ together with two natural
transformations:
\begin{itemize}
  \item \textbf{Unit} (return): $\eta \colon \mathrm{Id} \Rightarrow T$
  \item \textbf{Multiplication} (join/flatten):
        $\mu \colon T \circ T \Rightarrow T$
\end{itemize}
Subject to:
\begin{itemize}
  \item \textbf{Associativity}: $\mu \circ T\mu = \mu \circ \mu T$
        (flattening twice is the same regardless of order)
  \item \textbf{Unit laws}: $\mu \circ T\eta = \mu \circ \eta T = \mathrm{id}_T$
        (wrapping then flattening is a no-op)
\end{itemize}
Equivalently, via \textbf{bind} (Kleisli extension):
given $f \colon A \to T(B)$, define
$\mathrm{bind}(m, f) \colon T(A) \to T(B)$.
\end{formal}

Don't panic. Let's see the three monads hiding in Tech World.

\subsection{The Future Monad: \texttt{async}/\texttt{await}}

\begin{youknow}
Every \texttt{async} function in Dart returns a \texttt{Future<T>}.
\texttt{Future} is a monad:
\begin{itemize}
  \item \textbf{Unit}: \texttt{Future.value(x)} --- wrap a value in a completed future
  \item \textbf{Bind}: \texttt{future.then(f)} --- when the future completes,
        apply \texttt{f} and get a new future
  \item \texttt{await} is syntactic sugar for \texttt{.then}
\end{itemize}
\end{youknow}

In Tech World, the room connection sequence is monadic composition:

\begin{lstlisting}
// lib/rooms/room_session.dart
Future<void> connect() async {
  await _livekitService.connect(url, token);  // Future<void>
  await _livekitService.enableMedia();         // Future<void>
  _listenForConnectionLoss();                  // void (sync)
}
\end{lstlisting}

Each \texttt{await} is a bind: ``when the previous step completes,
do the next step.'' The monad laws guarantee that:
\begin{enumerate}
  \item Wrapping a value and immediately awaiting it gives back the value
        (unit law)
  \item The order of \texttt{await} grouping doesn't matter (associativity)
\end{enumerate}

\begin{reallife}
A monad is like a to-do list where each item might take time.
``Buy ingredients'' $\to$ ``Cook dinner'' $\to$ ``Wash up.''
You can't start cooking until shopping is done (that's bind).
``Do nothing and then cook'' is the same as ``cook'' (unit law).
And grouping items --- ``(buy, cook), wash'' vs ``buy, (cook, wash)'' ---
gives the same evening (associativity).
\end{reallife}


\subsection{The Maybe Monad: Null-Safe Chaining}

Dart's null-safety operator \texttt{?.} is the Maybe monad in disguise:

\begin{lstlisting}
_oracleService?.cancel();
currentRoom?.leave();
\end{lstlisting}

Here $T(A) = A\texttt{?}$ (nullable type). Unit is the identity (a non-null value),
and bind is the \texttt{?.} operator: if the left side is \texttt{null}, propagate
\texttt{null}; otherwise, apply the function.

The monad laws hold: chaining \texttt{?.} is associative, and starting with a
non-null value then immediately using \texttt{?.} does nothing.


\subsection{The Writer Monad: Accumulating Events}
\label{sec:writer}

This is the most interesting monad in Tech World --- and the developers
built it without knowing it was a monad.

\begin{formal}
The \textbf{Writer monad} for a monoid $(M, \cdot, e)$ is:
\[
W(A) = A \times M
\]
with:
\begin{itemize}
  \item \textbf{Unit}: $\eta(a) = (a, e)$ --- pair the value with the empty log
  \item \textbf{Bind}: $(a, m) \mathbin{>\!\!>\!\!=} f = \text{let } (b, m') = f(a) \text{ in } (b, m \cdot m')$
\end{itemize}
The monoid operation $\cdot$ accumulates the ``side output'' (a log, a list of
events, etc.).
\end{formal}

\begin{youknow}
Tech World defines:
\begin{lstlisting}
typedef WithEvents<T> = (T, List<AppEvent>);
\end{lstlisting}
This is exactly $W(T) = T \times \texttt{List<AppEvent>}$ where
the monoid is $(\texttt{List<AppEvent>}, \texttt{++}, \texttt{[]})$
--- list concatenation with empty list as identity.

Business functions return their result \emph{plus} a list of events
that should be dispatched:
\begin{lstlisting}
// lib/spellbook/cast_effects.dart
List<AppEvent> castSuccessEvents(PromptChallengeId id) {
  return [
    SpellCast(challengeId: id),
    ChallengeCompleted(ref: PromptRef(id)),
  ];
}
\end{lstlisting}

The caller collects events from multiple functions and dispatches them all:
\begin{lstlisting}
final events1 = castSuccessEvents(challengeId);
final events2 = unlockDoorEvents(doorId);
dispatch([...events1, ...events2]);  // monoid operation: list concatenation
\end{lstlisting}
The \texttt{[...events1, ...events2]} is the monoid multiplication $m \cdot m'$!
\end{youknow}

\subsubsection*{Why This Matters}

The Writer monad laws guarantee that events accumulate correctly:

\begin{itemize}
  \item \textbf{Left unit}: Starting with no events and then running a function
        gives the same events as just running the function.
        $\eta(a) \mathbin{>\!\!>\!\!=} f = f(a)$
  \item \textbf{Right unit}: Running a function and then ``adding no events''
        changes nothing.
        $m \mathbin{>\!\!>\!\!=} \eta = m$
  \item \textbf{Associativity}: Composing three event-producing functions
        gives the same total event list regardless of grouping.
        $(m \mathbin{>\!\!>\!\!=} f) \mathbin{>\!\!>\!\!=} g = m \mathbin{>\!\!>\!\!=} (\lambda x.\, f(x) \mathbin{>\!\!>\!\!=} g)$
\end{itemize}

The associativity law is the critical one. If it fails, composing three
business functions could \emph{lose events from the middle one}.
In production, that means incomplete audit logs, missed analytics, or
--- worse --- a spell cast that the server never hears about.

\begin{reallife}
A writer monad is like a group of friends writing a shared travel diary.
Each person (function) has experiences (a result) and writes diary entries
(events). At the end of the trip, you concatenate everyone's entries
into one diary. The monad laws say:
\begin{itemize}
  \item Someone who didn't go on the trip contributes an empty diary (unit).
  \item The order you staple diaries together doesn't matter (associativity).
  \item No entries get lost in the stapling process.
\end{itemize}
\end{reallife}


%=============================================================================
\section{Comonads: Consuming Context}
\label{sec:comonads}
%=============================================================================

If a monad is about \emph{producing} effects, a comonad is about
\emph{consuming} context. It's the dual notion.

\begin{formal}
A \textbf{comonad} on a category $\mathcal{C}$ is a functor
$W \colon \mathcal{C} \to \mathcal{C}$ together with:
\begin{itemize}
  \item \textbf{Counit} (extract): $\epsilon \colon W \Rightarrow \mathrm{Id}$
        --- extract the value from the context
  \item \textbf{Comultiplication} (duplicate):
        $\delta \colon W \Rightarrow W \circ W$
        --- every context contains a context-of-contexts
\end{itemize}
\end{formal}

\subsection{ValueNotifier: The Store Comonad (Partial)}

Flutter's \texttt{ValueNotifier<T>} is a reactive state cell. In Tech World:

\begin{lstlisting}
// lib/chat/chat_service.dart
final ValueNotifier<BotStatus> _botStatus =
    ValueNotifier(BotStatus.absent);

ValueListenable<BotStatus> get botStatus => _botStatus;
\end{lstlisting}

The comonad structure:
\begin{itemize}
  \item \textbf{Extract}: \texttt{notifier.value} --- get the current value
  \item \textbf{Extend} (missing): would be
        ``given a function that reads the notifier and produces a value,
        create a \emph{new} notifier that tracks the derived value.''
\end{itemize}

Flutter doesn't provide \texttt{extend}, but the \texttt{extract} half is
everywhere: widgets read \texttt{.value} to get the current state.

\begin{reallife}
A comonad is like a spreadsheet cell.
\textbf{Extract} reads the current value.
\textbf{Extend} is like writing a formula: ``this cell equals
\emph{that} cell times 2.'' The formula creates a new cell (context)
from an existing one.
Every spreadsheet user understands comonads --- they just call them
``formulas.''
\end{reallife}

\subsection{The Flame Component Tree: A Rose Tree Comonad}

The Flame game engine organises everything into a tree of \texttt{Component}
objects. This is a \textbf{rose tree} --- a tree where each node can have
any number of children.

\begin{lstlisting}
// Flame Component tree structure (simplified)
class Component {
  List<Component> children;
  Component? parent;
  void add(Component child) { ... }
  void remove(Component child) { ... }
}
\end{lstlisting}

A rose tree carries a comonad:
\begin{itemize}
  \item \textbf{Extract}: the data at the root node (this component)
  \item \textbf{Duplicate}: replace each node with the subtree rooted at
        that node --- every node ``sees'' its entire subtree
\end{itemize}

In Tech World, the \texttt{BubbleManager} receives an \texttt{addComponent}
callback rather than owning the tree directly:
\begin{lstlisting}
BubbleManager({
  required void Function(Component) addComponent,
})
\end{lstlisting}
This is a \textbf{co-Kleisli arrow}: a function $W(A) \to B$ that
consumes context (the component tree) to produce a result (adding a child).


%=============================================================================
\section{Kleisli Categories: Composing Effectful Functions}
\label{sec:kleisli}
%=============================================================================

\begin{formal}
Given a monad $T$ on $\mathcal{C}$, the \textbf{Kleisli category}
$\mathcal{C}_T$ has:
\begin{itemize}
  \item The same objects as $\mathcal{C}$
  \item Morphisms $A \to B$ in $\mathcal{C}_T$ are morphisms
        $A \to T(B)$ in $\mathcal{C}$ (``effectful functions'')
  \item Composition: $g \circ_T f = \mu \circ T(g) \circ f$
        (run $f$, then lift $g$ into the monad and flatten)
  \item Identity: $\eta_A$
\end{itemize}
\end{formal}

\begin{youknow}
Every chain of \texttt{async} functions is Kleisli composition in the
Future monad's Kleisli category:
\begin{lstlisting}
// Each function is a Kleisli arrow: A -> Future<B>
Future<Room> connect(String url, String token) async { ... }
Future<void> enableMedia(Room room) async { ... }
Future<void> publishState(Room room) async { ... }

// Kleisli composition via await:
final room = await connect(url, token);    // String -> Future<Room>
await enableMedia(room);                    // Room -> Future<void>
await publishState(room);                   // Room -> Future<void>
\end{lstlisting}
The \texttt{await} keyword is literally Kleisli composition:
``unwrap the future result, feed it to the next effectful function.''
\end{youknow}

\subsection{Writer Kleisli: Event-Producing Function Chains}

The Writer monad also has a Kleisli category. Tech World's business functions
are Kleisli arrows $A \to (B, \texttt{List<AppEvent>})$:

\begin{lstlisting}
// Kleisli arrow: PromptChallengeId -> (void, List<AppEvent>)
List<AppEvent> castSuccessEvents(PromptChallengeId id) { ... }

// Kleisli arrow: PromptChallengeId -> Future<List<AppEvent>>
Future<List<AppEvent>> applyCastSuccessEffects(...) async { ... }
\end{lstlisting}

Composition is done manually by concatenating event lists:
\texttt{[...events1, ...events2]}. This \emph{is} Kleisli composition,
just without an explicit \texttt{composeK} function.

\begin{reallife}
Imagine a relay race where each runner carries a bag of souvenirs.
Each runner (Kleisli arrow) runs their leg (computes a result) and
picks up souvenirs along the way (accumulates events).
At each handoff, the new runner receives the baton (the result)
and the bags are combined (events concatenated).
At the finish line, you have one big bag of all souvenirs
(the complete event list, dispatched at the end).
\end{reallife}


%=============================================================================
\section{Monoids: Things You Can Combine}
\label{sec:monoids}
%=============================================================================

\begin{formal}
A \textbf{monoid} is a set $M$ with a binary operation
$\cdot \colon M \times M \to M$ and an identity element $e \in M$ such that:
\begin{itemize}
  \item \textbf{Associativity}: $(a \cdot b) \cdot c = a \cdot (b \cdot c)$
  \item \textbf{Identity}: $e \cdot a = a \cdot e = a$
\end{itemize}
A monoid is a category with one object --- the morphisms are elements of $M$,
and composition is the monoid operation.
\end{formal}

\subsection{The Event List Monoid}

We already met this one: $(\texttt{List<AppEvent>}, \texttt{++}, \texttt{[]})$
is the \textbf{free monoid} on \texttt{AppEvent}. It powers the Writer monad
(Section~\ref{sec:writer}).

\subsection{The CRDT Join-Semilattice}
\label{sec:crdt}

Tech World has a collaborative map editor where multiple players can edit
the game world simultaneously. Edits are synchronised using a
\textbf{Conflict-free Replicated Data Type} (CRDT).

\begin{lstlisting}
// lib/map_editor/crdt/cell_version_map.dart
bool shouldApply(String cellKey, int counter, String playerId) {
  final existing = _versions[cellKey];
  if (existing == null) return true;
  // Compare (counter, playerId) lexicographically
  if (counter > existing.counter) return true;
  if (counter == existing.counter &&
      playerId.compareTo(existing.playerId) > 0) return true;
  return false;
}
\end{lstlisting}

The merge operation is: take the \emph{maximum} version for each cell,
where versions are ordered by \texttt{(counter, playerId)} lexicographically.

\begin{formal}
A \textbf{join-semilattice} is a partially ordered set where every pair of
elements has a least upper bound (join):
\[
a \vee b = \max(a, b)
\]
A join-semilattice is an idempotent commutative monoid:
\begin{itemize}
  \item \textbf{Associativity}: $(a \vee b) \vee c = a \vee (b \vee c)$
  \item \textbf{Commutativity}: $a \vee b = b \vee a$
  \item \textbf{Idempotence}: $a \vee a = a$
\end{itemize}
\end{formal}

All three laws hold for Tech World's CRDT merge:
\begin{itemize}
  \item \textbf{Associativity}: \texttt{max} over a total order is associative. \checkmark
  \item \textbf{Commutativity}: \texttt{max(a, b) = max(b, a)}. \checkmark
  \item \textbf{Idempotence}: \texttt{max(a, a) = a}. \checkmark
\end{itemize}

\begin{theorem}[CRDT Convergence]
If the merge operation forms a join-semilattice, then all replicas that
have received the same set of operations (in any order) converge to the
same state.
\end{theorem}
\begin{proof}[Proof sketch]
Since join is associative, commutative, and idempotent, the final state
depends only on the \emph{set} of operations, not their order or
multiplicity. Any permutation or repetition yields the same result.
\end{proof}

This is why the collaborative map editor works: two players making edits
in different orders will always converge. Category theory (specifically,
the algebraic structure of the merge) is the mathematical guarantee
backing this feature.

\begin{reallife}
Imagine two friends independently updating a shared grocery list.
The merge rule is ``keep the latest version of each item.''
Friend A writes ``3 apples'' at 2pm; Friend B writes ``5 apples'' at 3pm.
Regardless of the order the server receives these updates, the merged
list says ``5 apples'' (the later timestamp wins).
That's a join-semilattice: the latest version is the join.
Idempotence means receiving the same update twice doesn't change anything.
\end{reallife}


%=============================================================================
\section{Adjunctions: Perfect Partnerships}
\label{sec:adjunctions}
%=============================================================================

\begin{formal}
An \textbf{adjunction} $F \dashv G$ between categories $\mathcal{C}$ and
$\mathcal{D}$ consists of functors $F \colon \mathcal{C} \to \mathcal{D}$
and $G \colon \mathcal{D} \to \mathcal{C}$ such that:
\[
\mathrm{Hom}_{\mathcal{D}}(F(A), B) \cong \mathrm{Hom}_{\mathcal{C}}(A, G(B))
\]
naturally in $A$ and $B$. We say $F$ is the \textbf{left adjoint} and $G$
is the \textbf{right adjoint}.
\end{formal}

This is the most abstract definition so far, but the idea is concrete:
$F$ and $G$ are ``perfect partners'' --- anything you can do with $F$
on one side, you can equivalently do with $G$ on the other.

\subsection{The Serialisation Adjunction}

Tech World's \texttt{toJson}/\texttt{fromJson} pair forms an adjunction
between Dart types and JSON structures:

\begin{lstlisting}
// Left adjoint: Dart -> JSON
Map<String, dynamic> toJson();

// Right adjoint: JSON -> Dart
factory MapEditBatch.fromJson(Map<String, dynamic> json);
\end{lstlisting}

The adjunction says: a map from ``JSON representation of a Dart object''
to some target is the same as a map from the Dart object to the ``Dart
interpretation of that target.'' In practice, this means:

\begin{itemize}
  \item \textbf{Round-trip} (unit of adjunction):
        \texttt{fromJson(toJson(x)) == x}
  \item \textbf{Retraction} (counit of adjunction):
        \texttt{toJson(fromJson(j))} produces a ``canonical'' JSON
\end{itemize}

The round-trip property is what makes the LiveKit data channel work:
one client serialises a map edit batch to JSON, sends it over the wire,
and another client deserialises it back to the same Dart object.

\begin{reallife}
Translation between languages is an adjunction (approximately).
Translating a book to French (\texttt{toJson}) and back to English
(\texttt{fromJson}) should give you the original book (round-trip).
In practice, something is always lost in translation --- just as
Dart's type information is partially lost in \texttt{dynamic} JSON.
The closer the adjunction is to an isomorphism,
the more faithful the translation.
\end{reallife}

\subsection{The Event--Sink Adjunction}

The deepest adjunction in Tech World connects the event-producing side
(Writer monad) with the event-consuming side (dispatch/sinks):

\[
\underbrace{\text{Produce events}}_{\text{Left adjoint } F}
\quad \dashv \quad
\underbrace{\text{Consume via sinks}}_{\text{Right adjoint } G}
\]

\begin{itemize}
  \item \textbf{Left adjoint} ($F$): Business functions that return
        \texttt{(result, List<AppEvent>)} --- they \emph{produce} events.
  \item \textbf{Right adjoint} ($G$): Sinks registered with the dispatch
        system --- they \emph{consume} events.
  \item \textbf{Counit}: The \texttt{dispatch} function itself ---
        it connects production to consumption, sending accumulated events
        to all registered sinks including LiveKit (\texttt{publishJson}).
\end{itemize}

The adjunction tells us these two sides \emph{must} be designed together.
A new event type requires both a producer (left) and handling in every
sink (right). The sealed class on \texttt{AppEvent} enforces this at
compile time.


%=============================================================================
\section{The Yoneda Perspective: Build Up, Apply Once}
\label{sec:yoneda}
%=============================================================================

\begin{formal}
The \textbf{Yoneda lemma} states: for a functor $F \colon \mathcal{C} \to \mathbf{Set}$
and an object $A$, there is a natural bijection:
\[
\mathrm{Nat}(\mathrm{Hom}(A, -), F) \cong F(A)
\]
A natural transformation from the hom-functor at $A$ to $F$ is completely
determined by a single element of $F(A)$.
\end{formal}

The Yoneda lemma is deep and general, but it has a simple computational
consequence:

\begin{youknow}
\textbf{Map fusion}: instead of mapping over a structure three times:
\begin{lstlisting}
list.map(f).map(g).map(h)  // three traversals
\end{lstlisting}
you can fuse into one traversal:
\begin{lstlisting}
list.map((x) => h(g(f(x))))  // one traversal
\end{lstlisting}
The Yoneda lemma guarantees these are equal.
More generally: \textbf{build up the transformation, apply it once.}
\end{youknow}

In Tech World, the Stream pipeline for LiveKit data does exactly this ---
multiple \texttt{.map} and \texttt{.where} calls are chained on a stream.
Dart's Stream is lazy, so the transformations are already fused into a single
pass per element. This is the Yoneda optimisation happening automatically.

\begin{reallife}
Yoneda is like a restaurant order.
Instead of: ``bring me bread'' (waiter walks to kitchen),
``add butter'' (waiter walks again),
``also some jam'' (waiter walks a third time) ---
you say: ``bring me bread with butter and jam'' (one trip).
Building up the complete request before executing it is the Yoneda
principle. Apply the composed function once, not each piece separately.
\end{reallife}


%=============================================================================
\section{Summary: A Categorical Map of Tech World}
\label{sec:summary}
%=============================================================================

Here is every categorical structure found in the codebase, in one table:

\begin{center}
\small
\begin{tabular}{@{}llll@{}}
\toprule
\textbf{Structure} & \textbf{Category Theory} & \textbf{Where in Code} & \textbf{Real-Life Analogy} \\
\midrule
Dart types + functions & Category $\mathbf{Dart}$ & Everywhere & City road network \\
\texttt{sealed class} & Coproduct & \texttt{AppEvent}, \texttt{ChallengeRef} & Hospital triage \\
\texttt{PiiPolicy} enum & 2-element poset category & \texttt{events/pii\_policy.dart} & Security clearance \\
\texttt{Stream.map} & Endofunctor & LiveKit data pipeline & Assembly line \\
\texttt{Direction} enum & Functor to $\mathbb{R}^2$ & \texttt{flame/shared/direction.dart} & Dictionary \\
\texttt{wireName/parse} & Natural transformation & All domain enums & Language translation \\
\texttt{toJson/fromJson} & Adjunction unit/counit & Events, CRDT ops & Book translation \\
Logger bridge & Partial nat.\ transformation & \texttt{events/logger\_bridge.dart} & Censored translation \\
\texttt{Future} & Monad & Everywhere & To-do list \\
\texttt{?.} null-safe & Maybe monad & Service calls & ``If present, then\ldots'' \\
\texttt{WithEvents<T>} & Writer monad & \texttt{events/types.dart} & Shared travel diary \\
\texttt{ValueNotifier} & Store comonad (partial) & Chat, room session & Spreadsheet cell \\
Component tree & Rose tree comonad & Flame engine & Family tree \\
\texttt{async/await} chains & Future Kleisli category & Room connection & Relay race \\
Event $\to$ sink & Adjunction & Dispatch system & Author $\leftrightarrow$ reader \\
CRDT merge & Join-semilattice & \texttt{crdt/cell\_version\_map.dart} & Shared grocery list \\
\texttt{List<AppEvent>} & Free monoid & Event accumulation & Diary entries \\
Lamport clock & Logical time & \texttt{crdt/undo\_manager.dart} & Postal timestamps \\
Stream pipeline & Yoneda fusion & LiveKit data channel & Restaurant order \\
\bottomrule
\end{tabular}
\end{center}

\section{What Category Theory Cannot Tell You}

Intellectual honesty requires us to say where category theory
\emph{doesn't} help:

\begin{enumerate}
  \item \textbf{Network reliability.} Reconnection backoff timing
        (2s/4s/8s) is an engineering judgment, not a categorical one.
  \item \textbf{Game design.} Whether the proximity threshold should be
        3 grid squares or 5 is a UX decision.
  \item \textbf{Privacy classification.} Whether a field is PII is a legal
        question. Category theory ensures the gate works once you've decided.
  \item \textbf{AI quality.} The oracle service wraps an LLM. The quality
        of its responses is not a mathematical property.
  \item \textbf{Render performance.} Frame rates, shader compilation,
        physics simulation --- these are engineering concerns.
  \item \textbf{Database schema.} Whether to use arrays or subcollections
        in Firestore is a cost/latency tradeoff.
  \item \textbf{Global state.} Whether the \texttt{Locator} singleton is
        good architecture is a design opinion, not a theorem.
\end{enumerate}

Category theory is a lens, not a hammer. It reveals structure where structure
exists. Applying it where it doesn't belong produces only noise.

\vspace{2em}
\begin{center}
\rule{0.5\textwidth}{0.4pt}\\[1em]
\textit{``Category theory is not about categories any more than\\
number theory is about numbers.'' --- attributed to various}
\end{center}

\end{document}
